% Lösungen Einheit 4 von 4 – Abstand als Argument (zusammenbau v0.1)
\einheitenkopf{Einheit 4 von 4 · Abstand als Argument}

\erg{16}{a) senkrecht: die Strecke ist der Abstand \quad b) größer: $\overrightarrow{BE} = (0 | 0 | -3)$ ist kein Vielfaches des Normalenvektors (1 | 2 | 2), BE steht nicht senkrecht auf den Ebenen, und das Lot ist die kürzeste Verbindung (Kontrolle: $|\overrightarrow{BE}| = 3$, Ebenenabstand $= 2$) \quad c) Das Lot von F auf t trifft t im Lotfußpunkt L(2,4 | 1,2) links von G, außerhalb der Strecke; nur L hat von F den Abstand 3,1, jeder Punkt von GJ ist weiter entfernt. \quad d) Die Tafel steht senkrecht, jeder ihrer Punkte hat denselben waagerechten Abstand zur Achse wie der Punkt der Oberkante darüber; auf DE ist M wegen der Symmetrie von D und E der Punkt mit dem kleinsten Abstand zur Achse – also ist der Abstand der Tafel der Abstand von M. \quad e) Oberkante der Öffnung: $3y + 2 \cdot 1{,}2 = 15$, $y = 4{,}2$; Länge in der Wandebene $\sqrt{(5 - 4{,}2)^2 + 1{,}2^2} \approx 1{,}44$ m (nicht der waagerechte Abstand 0,8 m) \quad f) Wegen der Symmetrie liegt die Lampe auf der Achse MS: L(5 | 3 | z); $\tfrac{|12 - 3z|}{5} = 1$ ergibt $z = \tfrac{7}{3}$ oder $z = \tfrac{17}{3}$ (über der Spitze, außerhalb), also L(5 | 3 | 2,33) \quad g) I: Lotgerade durch P mit dem Normalenvektor (0 | 2 | 1) von F; II: Q liegt in F, ist also der Lotfußpunkt; III: der Abstand zu F ist gleich dem Abstand p zur Grundfläche – P ist der Mittelpunkt der Kugel, die alle Flächen berührt.}

16c)

\begin{ksys}[xmin=0,xmax=8,ymin=0,ymax=8,klein]
\gerade{0.5}{0}{t}
\punkt{4}{2}{G}
\punkt{8}{4}{J}
\punkt{1}{4}{F}
\gerade{-2}{6}{}
\punkt{2.4}{1.2}{L}
\end{ksys}

\erg{17}{a) Fehler: die schräge Strecke AB als Ebenenabstand genommen; $\overrightarrow{AB} = (1 | -1 | -5)$ ist kein Vielfaches von (2 | 1 | 2). Richtig: $d = \tfrac{|11 - 2|}{3} = 3$ \quad b) Jede andere Verbindung zu einem Ebenenpunkt bildet mit dem Lot und der Strecke zwischen den Fußpunkten ein rechtwinkliges Dreieck, in dem sie die Hypotenuse ist; die Hypotenuse ist länger als jede Kathete. \quad c) innere Teilung: $P_1 = A + \tfrac{2}{3}\overrightarrow{AB} = (6 | 3 | 0)$; äußere Teilung: $P_2 = A + 2\overrightarrow{AB} = (14 | 7 | 0)$}
